\section{3. Funcionamiento general de EBP}
EBP actúa como punto central entre las fuentes de contenido y los receptores. Una señal es recibida por la plataforma, administrada como un flujo independiente y posteriormente puesta a disposición de uno o varios destinos compatibles. La supervisión del proceso se realiza mediante el NOC.

\begin{tabularx}{\textwidth}{|Y|Y|Y|}\hline
\rowcolor{Navy}\THead{FUENTES} & \THead{EBP} & \THead{DESTINOS} \\ \hline
ENLACE\\EJTV\\IMPACT\\Otros proveedores & \textbf{Recepción}\\\textbf{Administración}\\\textbf{Distribución}\\\textbf{Supervisión} & Cable Centro\\Otros operadores\\Clientes\\Aplicaciones web\\Nuevos destinos \\ \hline
\end{tabularx}

En este contexto, \textbf{Reader} identifica al receptor o consumidor de una señal distribuida por EBP. Puede corresponder a un operador, un sistema receptor, un cliente o una aplicación. Los protocolos son los mecanismos técnicos utilizados para transportar o entregar el contenido; actualmente la plataforma trabaja con SRT, RTMP, RTSP, HLS y WebRTC.

La arquitectura no está limitada a contenido propio. Pueden incorporarse señales suministradas por otros proveedores, siempre que exista compatibilidad técnica y la correspondiente autorización para su uso y distribución.
